\label{sec:mixing}

This section solves the allocation problem in two steps. First, \Cref{ssec:mixing} fixes the calibration budget $x$ and chooses the best mixing weight $\alpha$ between the real and synthetic estimators. This gives a one-dimensional risk curve $R_n(x)$. Second, \Cref{ssec:foc} differentiates that curve with respect to $x$ and gives the corrected first-order condition. The difference from the usual heuristic is one retained term, the synthetic sampling variance $v_n$, and that term drives the phase diagram in \Cref{sec:phase}. Finally, \Cref{ssec:safe} gives a finite-sample safety check for falling back to all-real estimation. All proofs are deferred to Appendix~\ref{app:proofs}.

\subsection{Optimal mixing for fixed allocation}\label{ssec:mixing}

We first compute the leading-order MSE of the linear allocation estimator and minimize it over the mixing weight. Throughout, $\hat\theta^{\mathrm{L}}(\alpha,x)$ is as in \cref{eq:linest}, $V_R(x)=a/(n-x)$, and $B_n(x)=v_n+cx^{-2\beta}$.

\begin{lemma}[Working MSE]\label{lem:mse-decomp}
  Under Assumptions~\ref{asm:al-real}, ~\ref{asm:al-synth} and ~\ref{asm:pl},
  \begin{equation*}
    \mathbb{E}\bigl[(\hat\theta^{\mathrm{L}}(\alpha,x)-\theta^*)^2\bigr]
    \;=\;\alpha^2\,V_R(x)\;+\;(1-\alpha)^2\,B_n(x)\;+\;o\bigl(V_R(x)\bigr)+o\bigl(B_n(x)\bigr),
  \end{equation*}
  uniformly in $\alpha\in[0,1], x\in (0,n)$.
\end{lemma}

In words, sample-splitting makes $\hat\theta_R(E_x)$ and $\hat\theta_S(x)$ unconditionally independent, so the cross term in the squared error vanishes and the MSE decomposes additively. The first term is the variance contribution from the real estimator, downweighted by $\alpha^2$; the second is the unconditional MSE of the synthetic estimator around $\theta^*$, downweighted by $(1-\alpha)^2$.

Minimizing the leading-order MSE in $\alpha$ for fixed $x$ is a one-dimensional convex problem and admits a closed-form solution.

\begin{proposition}[Optimal mixing weight and profiled risk]\label{prop:mix}
  Fix $x\in(0,n)$ such that $0<B_n(x)<\infty$. Then the unique minimizer of $\alpha\mapsto \overline{\mathrm{MSE}}_n(\alpha,x)$ over $\alpha\in[0,1]$, where $\overline{\mathrm{MSE}}_n(\alpha,x) := \alpha^2 V_R(x)+(1-\alpha)^2 B_n(x),$ is
  \begin{equation}\label{eq:alpha-star}
    \alpha^*(x)
    =
    \frac{B_n(x)}{V_R(x)+B_n(x)}
    \in(0,1).
  \end{equation}
  The corresponding profiled risk is
  \begin{equation}\label{eq:Rn}
    R_n(x)
    :=
    \overline{\mathrm{MSE}}_n(\alpha^*(x),x)
    =
    \frac{V_R(x)B_n(x)}{V_R(x)+B_n(x)}
    =
    \frac{aB_n(x)}{a+(n-x)B_n(x)}.
  \end{equation}
\end{proposition}

The formula has a simple reading. If $B_n(x)$ is large relative to $V_R(x)$, the synthetic estimator is unreliable, so the optimal rule puts most weight on real data ($\alpha^*\to 1$). If $B_n(x)$ is small, the synthetic estimator is accurate, so it gets most of the weight ($\alpha^*\to 0$). The profiled risk $R_n(x)$ is the usual inverse-variance pooling risk for two independent estimators, with $B_n(x)$ playing the role of the synthetic estimator's effective error. The remaining problem is to choose the calibration budget $x$ that minimizes this profiled risk.

\subsection{Exact corrected first-order condition}\label{ssec:foc}

The fixed-share shortcut used as a baseline treats the synthetic estimator as if its only error were calibration bias. In our notation, it differentiates a risk surrogate after replacing $B_n(x)=v_n+cx^{-2\beta}$ by $cx^{-2\beta}$. Dropping $v_n$ is harmless only when synthetic sampling variance is negligible at the relevant allocation. It is not harmless when the synthetic budget $m_n$ grows with $n$, because $v_n$ then enters the marginal value of one more calibration observation. Keeping this term produces a different first-order condition, one that leads to the phase transition in \Cref{sec:phase} rather than a rate-independent fixed-share rule such as $\lambda=2\beta/(1+2\beta)$.

\begin{proposition}[Exact first-order condition]\label{prop:foc}
  Under Assumption~\ref{asm:pl} and~\ref{asm:vn} with $v_n>0$, every interior critical point $x\in(0,n)$ of $R_n$ satisfies
  \begin{equation}\label{eq:foc}
    \bigl(v_n+c\,x^{-2\beta}\bigr)^2\;=\;2\beta\,a\,c\,x^{-(2\beta+1)},
  \end{equation}
  or equivalently, after multiplying through by $x^{2\beta+1}$,
  \begin{equation}\label{eq:foc-poly}
    v_n^2\,x^{2\beta+1}\;+\;2c\,v_n\,x\;+\;c^2\,x^{1-2\beta}\;=\;2\beta\,a\,c.
  \end{equation}
  For $\beta\le 1/2$ this equation has at most one solution on $(0,\infty)$; for $\beta>1/2$ it has at most two, and the active interior root is the larger.
\end{proposition}

The left-hand side of \cref{eq:foc} is $B_n(x)^2$, the squared effective synthetic error. The right-hand side is the cost, in real-estimator precision, of spending one more real observation on calibration. Thus the optimum balances synthetic-error reduction against real-data loss. The mixed term $2cv_n\,x$ in \cref{eq:foc-poly} is the piece lost when one drops synthetic sampling variance; keeping it is the algebraic content of the correction. The root structure is controlled by the auxiliary function $\Phi(x):=B_n(x)^2 x^{2\beta+1}$, which is monotone for $\beta\le 1/2$ and U-shaped for $\beta>1/2$. \Cref{sec:phase} turns that shape into the allocation phase diagram.

\subsection{Safe-improvement condition}\label{ssec:safe}

The corrected FOC characterizes interior optima of $R_n$, but the global oracle is taken over the closed interval $[0,n]$, so the boundary $x=0$ (no calibration, all-real estimator on $n$ observations, with risk $a/n$) must be compared against any interior candidate. The next result reduces this comparison to a one-line algebraic check.

\begin{proposition}[Exact safe-improvement condition]\label{prop:safe}
  For every $x\in(0,n)$,
  \begin{equation}\label{eq:safe}
    R_n(x)\;<\;\frac{a}{n}\quad\Longleftrightarrow\quad x\,B_n(x)\;<\;a.
  \end{equation}
  Equivalently, $x\,(v_n+cx^{-2\beta})<a$. The boundary risk is $R_n(0)=a/n$.
\end{proposition}

In words, calibration budget $x$ beats all-real estimation on $n$ observations exactly when $x\,B_n(x)<a$. The product $x\,B_n(x)$ compares the amount of real data spent on calibration with the remaining error of the calibrated synthetic estimator. If the inequality fails because $v_n$ is too large or because the bias term $cx^{-2\beta}$ would require too much calibration to shrink, the procedure returns the all-real estimator, $\hat x=0$, $\hat\alpha=1$. \Cref{ssec:foc} produces candidate interior optima; \cref{eq:safe} decides whether any candidate should be used. The adaptive procedure in \Cref{sec:adaptive} minimizes $R_n(x)$ on a grid that includes $x=0$ and $x=n$, then applies this safety check.

% =====================================================================
% POLISHED
